% !TeX root = ../../Book4.tex
% 纲要标题：### F.4 无穷范畴初步
% 纲要细目（写作范围）：
% #### F.4.1 单纯集与范畴的神经
% #### F.4.2 内角填充与映射空间
% ---

\section{无穷范畴初步}
\label{b4:sec:F04}

复形映射之间可以有同伦，同伦之间又可以有更高的比较。\subsecref{b4:subsec:2004:02}指出，只取同伦类会丢失这些资料。单纯集提供了一种组织方法：顶点记录对象，边记录态射，三角形记录复合的比较，更高维单纯形记录比较之间的相容性。普通范畴也能用同一语言表示，因此可以从熟悉的严格复合出发，看清无穷范畴增加了什么。

\subsection{单纯集与范畴的神经}
\label{b4:subsec:F04:01}

\begin{definition}\label{b4:def:F-simplicial-set}
对整数 \(n\ge0\)，令 \([n]=\{0<1<\cdots<n\}\)。以这些有限非空线序集为对象、以不减映射为态射，组成范畴 \(\Delta\)。函子
\[
 S:\Delta^{\mathrm{op}}\longrightarrow\mathbf{Set}
\]
称为\term[dan-chun-ji]{单纯集}[simplicial set]，记 \(S_n=S([n])\)，其元素称为 \(n\) 维单纯形。对不减映射 \(\alpha:[m]\rightarrow[n]\)，记 \(\alpha^*=S(\alpha):S_n\rightarrow S_m\)。单纯集之间的映射是这种函子之间的自然变换。
\end{definition}

这里 \([0]\) 只有一个元素，\([1]\) 有两个有序元素。对 \(n\ge1\)、\(0\le i\le n\)，删去 \([n]\) 的第 \(i\) 个元素给出单射 \(\delta^i:[n-1]\rightarrow[n]\)；对 \(n\ge0\)、\(0\le i\le n\)，把 \([n+1]\) 中相邻的 \(i,i+1\) 合并给出满射 \(\sigma^i:[n+1]\rightarrow[n]\)。相应的
\[
 \mathrm d_i=(\delta^i)^*:S_n\longrightarrow S_{n-1},
 \qquad
 s_i=(\sigma^i)^*:S_n\longrightarrow S_{n+1}
\]
分别称为\term[mian-yingshe]{面映射}[face map]与\term[tuihua-yingshe]{退化映射}[degeneracy map]。前者取出一个面，后者重复一个顶点；它们之间的恒等式来自 \(\Delta\) 中的复合及函子律。例如 \(i<j\) 时有 \(\mathrm d_i\mathrm d_j=\mathrm d_{j-1}\mathrm d_i\)。

\begin{definition}\label{b4:def:F-standard-simplex}
设 \(n\ge0\)。单纯集
\[
 (\Delta^n)_m=\operatorname{Hom}_{\Delta}([m],[n]),
 \qquad
 \alpha^*(\beta)=\beta\alpha
\]
称为\term[biaozhun-danchunxing]{标准单纯形}[standard simplex]。其中 \(\alpha:[\ell]\rightarrow[m]\)、\(\beta:[m]\rightarrow[n]\) 为不减映射。
\end{definition}

标准单纯形的顶点为 \(0,\ldots,n\)，一个 \(m\) 维单纯形是一串允许重复的顶点 \(a_0\le\cdots\le a_m\)。不重复的串给出通常的面，重复的串是退化单纯形。对任意单纯集 \(S\)，给出 \(S_n\) 的一个元素 \(u\)，等价于给出映射 \(\Delta^n\rightarrow S\)：它在 \(\beta:[m]\rightarrow[n]\) 上的值必须为 \(\beta^*u\)，而在 \(1_{[n]}\) 上取回 \(u\)。

\begin{definition}\label{b4:def:F-category-nerve}
设 \(\mathcal A\) 为小范畴。把 \([n]\) 看作范畴：对象为 \(0,\ldots,n\)，当且仅当 \(a\le b\) 时有唯一的态射 \(a\rightarrow b\)。令
\[
 N(\mathcal A)_n=\operatorname{Ob}\operatorname{Fun}([n],\mathcal A),
 \qquad
 \alpha^*(F)=F\alpha,
\]
其中右边是函子范畴 \(\operatorname{Fun}([n],\mathcal A)\) 的对象集合。这个单纯集称为 \(\mathcal A\) 的\term[fanchou-shenjing]{神经}[nerve]。
\end{definition}

一个这样的函子由可复合的态射串
\[
 A_0\xrightarrow{f_1}A_1\xrightarrow{f_2}\cdots
 \xrightarrow{f_n}A_n
\]
唯一决定：从 \(a\) 到 \(b>a\) 的态射送到 \(f_b\cdots f_{a+1}\)。面映射 \(\mathrm d_0\) 删去 \(A_0,f_1\)，\(\mathrm d_n\) 删去 \(f_n,A_n\)；对 \(0<i<n\)，\(\mathrm d_i\) 删去 \(A_i\)，并用 \(f_{i+1}f_i\) 替代相邻的两个态射。退化映射 \(s_i\) 则重复 \(A_i\)，在两个副本之间插入 \(1_{A_i}\)。这些公式也说明 \(N(\mathcal A)_0\) 记录对象，\(N(\mathcal A)_1\) 记录态射，而 \(N(\mathcal A)_2\) 记录复合。

\ProseParagraphBreak
例如 \(N([1])=\Delta^1\)：从 \([n]\) 到 \([1]\) 的函子就是不减映射。它有一条从 \(0\) 到 \(1\) 的非退化边，却没有从 \(1\) 到 \(0\) 的边；单纯集中的边并不因能够组成三角形就必须可逆。

\subsection{内角填充与映射空间}
\label{b4:subsec:F04:02}

一个二维三角形若已给出两条首尾相接的边，补出整个三角形便同时给出了第三条边和复合的比较。高维情形要补的是已经相容的一组面，这由角的填充条件表达。

\begin{definition}\label{b4:def:F-horn-quasicategory}
设 \(n\ge1\)、\(0\le i\le n\)。以 \(\partial_j\Delta^n\) 表示由跳过顶点 \(j\) 的单射 \(\delta^j\) 诱导的 \(\Delta^{n-1}\) 在 \(\Delta^n\) 中的像。子单纯集
\[
 \Lambda_i^n=\bigcup_{j\ne i}\partial_j\Delta^n
 \ \subseteq\ \Delta^n
\]
称为第 \(i\) 个\term[jiao-danchunji]{角}[horn]；它包含除第 \(i\) 个面以外的全部余维一面及这些面的面。对 \(0<i<n\)，称它为内角。若单纯集 \(\mathcal C\) 满足：对每个 \(n\ge2\)、每个 \(0<i<n\) 和每个单纯集映射 \(u:\Lambda_i^n\rightarrow\mathcal C\)，都存在 \(\bar u:\Delta^n\rightarrow\mathcal C\)，使 \(\bar u|_{\Lambda_i^n}=u\)，则称 \(\mathcal C\) 为\term[nifanchou]{拟范畴}[quasi-category]。这是本附录采用的\term[wuqiongfanchou]{无穷范畴}[infinity-category]模型，即 \((\infty,1)\)-范畴模型。
\end{definition}

此定义要求填充存在，并不要求唯一，见 Lurie《Higher Topos Theory》作者2012年校订稿\cite[Definition 1.1.2.4]{Lurie2009HigherToposTheory}。二维内角 \(\Lambda_1^2\) 由边 \(0\rightarrow1\)、\(1\rightarrow2\) 组成；填充给出边 \(0\rightarrow2\) 及其与前两条边复合的比较。三维填充记录两种连续复合的相容性。普通范畴的严格复合使这些填充有更强的性质。

\ProseParagraphBreak
下图左边把二维角中已有的两条边画成实线，缺失的边画成虚线；还须填入整个三角形。右边是 \(\Lambda_1^3\)：已有面为 \(012,013,123\)，缺失面为与顶点 \(1\) 相对的 \(023\)，以浅灰标出。三维角已经含全部六条边，缺的是这个面以及整个三维单纯形，不能把它理解为少了一条边。
\[
\begin{tikzpicture}[scale=0.85]
 \coordinate (A) at (0,0);\coordinate (B) at (1.8,1.8);\coordinate (C) at (3.6,0);
 \draw[->] (A)--(B);\draw[->] (B)--(C);\draw[->,dashed] (A)--(C);
 \node[left] at (A) {\(0\)};\node[above] at (B) {\(1\)};\node[right] at (C) {\(2\)};
 \node at (1.8,-0.6) {\(\Lambda_1^2\)};
 \begin{scope}[xshift=6cm]
  \coordinate (D) at (0,0);\coordinate (E) at (1.4,0.55);\coordinate (F) at (3.5,0);\coordinate (G) at (1.7,2.4);
  \fill[gray!18] (D)--(F)--(G)--cycle;
  \draw[->] (D)--(E);\draw[->] (E)--(F);\draw[->] (E)--(G);
  \draw[->] (D)--(F);\draw[->] (D)--(G);\draw[->] (F)--(G);
  \node[left] at (D) {\(0\)};\node[below] at (E) {\(1\)};\node[right] at (F) {\(2\)};\node[above] at (G) {\(3\)};
  \node at (1.75,1.45) {缺面 \(023\)};
  \node at (1.75,-0.6) {\(\Lambda_1^3\)};
 \end{scope}
\end{tikzpicture}
\]

\begin{proposition}\label{b4:prop:F-nerve-inner-horns}
设 \(\mathcal A\) 为小范畴。对所有 \(n\ge2\)、\(0<i<n\)，每个映射 \(\Lambda_i^n\rightarrow N(\mathcal A)\) 都有唯一的填充 \(\Delta^n\rightarrow N(\mathcal A)\)。因而普通范畴的神经是拟范畴。
\end{proposition}
\begin{proof}
当 \(n=2\) 时，给出的角就是 \(A_0\xrightarrow{f_{01}}A_1\xrightarrow{f_{12}}A_2\)。唯一的填充对应这串态射的函子 \([2]\rightarrow\mathcal A\)，第三条边为 \(f_{02}=f_{12}f_{01}\)。

当 \(n=3\) 时，每条边都在保留的面中，记相应态射为 \(f_{ab}:A_a\rightarrow A_b\)。若 \(i=1\)，已给出三角形 \(012,013,123\)，故
\[
 f_{02}=f_{12}f_{01},\qquad
 f_{03}=f_{13}f_{01},\qquad
 f_{13}=f_{23}f_{12}.
\]
结合律给出 \(f_{03}=f_{23}(f_{12}f_{01})=f_{23}f_{02}\)，正是缺失面 \(023\) 所需的关系。若 \(i=2\)，已有面 \(012,023,123\)，则
\[
 f_{03}=f_{23}f_{02}
 =f_{23}(f_{12}f_{01})
 =(f_{23}f_{12})f_{01}=f_{13}f_{01},
\]
补出面 \(013\)。两种情形中，这些边均定义同一个函子 \([3]\rightarrow\mathcal A\)，且其限制与给出的三个面一致。

设 \(n\ge4\)。任取三个不同顶点 \(a<b<c\)，其余顶点至少有两个，故可以选出 \(j\notin\{a,b,c\}\) 且 \(j\ne i\)。于是三角形 \(abc\) 包含在保留的面 \(\partial_j\Delta^n\) 中。所有边也在角中，且它们在面的交上相容；记为 \(f_{ab}\)。由每个三角形属于范畴神经，得
\[
 f_{ac}=f_{bc}f_{ab}\qquad(a<b<c).
\]
将对象 \(a\) 送到 \(A_a\)，将 \(a\rightarrow b\) 送到 \(f_{ab}\)，将恒等态射送到恒等态射，便定义函子 \([n]\rightarrow\mathcal A\)。其与各保留面上的函子有相同的对象和边，因此限制相同，给出所求填充。对于 \(n\ge3\)，角已经包含全部边，而范畴神经的一个单纯形由这些边决定，故填充唯一；二维唯一性已在开头得到。
\end{proof}

这一命题是\cite[Proposition 1.1.2.2]{Lurie2009HigherToposTheory}中神经刻画的一方向。拟范畴允许同一可复合边对有不同的填充，这些选择由进一步的单纯形比较。\((\infty,1)\) 中的“\(1\)”表示一阶态射可以不可逆，而二阶及以上的比较在同伦意义下可逆；其精确含义由下面的映射空间体现。

\begin{definition}\label{b4:def:F-right-mapping-space}
设 \(\mathcal C\) 为拟范畴，\(x,y\in\mathcal C_0\) 为顶点。令
\[
 \operatorname{Map}_{\mathcal C}^{\mathrm R}(x,y)_n
 =\left\{\,u:\Delta^{n+1}\rightarrow\mathcal C
 \;\middle|\;
 u|_{\Delta^{\{0,\ldots,n\}}}\text{ 恒为 }x,
 \ u(n+1)=y\,\right\}.
\]
这里 \(\Delta^{\{0,\ldots,n\}}\) 是前 \(n+1\) 个顶点张成的面，“恒为 \(x\)”指该限制经 \(\Delta^0\) 分解。对 \(\alpha:[m]\rightarrow[n]\)，将它延拓为 \(\alpha^+:[m+1]\rightarrow[n+1]\)，在前 \(m+1\) 个顶点上取 \(\alpha\)，并令 \(\alpha^+(m+1)=n+1\)；以 \(u\mapsto u\alpha^+\) 为拉回。这给出一个单纯集，称为从 \(x\) 到 \(y\) 的\term[you-yingshekongjian]{右映射空间}[right mapping space]模型。
\end{definition}

其零维单纯形是边 \(x\rightarrow y\)。一维单纯形是一个顶点为 \(x,x,y\) 的三角形，底边为 \(1_x\)，另两边为两个待比较的态射。于是映射空间中的路径就是态射之间的同伦，更高单纯形继续记录这些同伦的比较。

\[
\begin{tikzcd}[column sep=4em,row sep=2.6em]
&y\arrow[d,phantom,"u" description]&\\
x\arrow[ur,"g"]\arrow[rr,"1_x"']&\phantom{x}&x\arrow[ul,"f"']
\end{tikzcd}
\]
这里底边是退化边 \(s_0x\)，两个底部顶点虽然同为对象 \(x\)，却是单纯形的第 \(0,1\) 个顶点；顶部是第 \(2\) 个顶点。若记 \(f=\mathrm d_0u\)、\(g=\mathrm d_1u\)，则 \(u\) 在右映射空间中就是端点为 \(f,g\) 的一维单纯形，而不是又一条 \(x\to y\) 的1态射。

\ProseParagraphBreak
填充条件使这一单纯集具有空间的性质。若单纯集 \(K\) 对所有 \(n\ge1\) 及 \(0\le i\le n\) 的角均有填充，则称它为\term[Kan-fuxing]{Kan 复形}[Kan complex]。
\footnote{康（Daniel Marinus Kan，1927--2013），荷兰数学家，主要研究同伦论，发展了单纯同伦方法与伴随函子理论。}
与拟范畴相比，这里还要求两个外角可填充。对拟范畴 \(\mathcal C\)，\(\operatorname{Map}_{\mathcal C}^{\mathrm R}(x,y)\) 是 Kan 复形，见\cite[Proposition 1.2.2.3]{Lurie2009HigherToposTheory}。该结果利用内角条件先得到右映射空间的一侧角填充，再由\cite[Proposition 1.2.5.1]{Lurie2009HigherToposTheory}补足另一侧；后一转换的角组合论证可在该处核读。以下将此模型记为 \(\operatorname{Map}_{\mathcal C}(x,y)\)。

\ProseParagraphBreak
空间中的路径在同伦意义下可逆，同伦之间的比较亦然；这解释了二阶及以上态射的可逆性。它没有要求 \(x\rightarrow y\) 在 \(\mathcal C\) 中可逆。例如对普通范畴 \(\mathcal A\)，右映射空间的每个 \(n\) 维单纯形都由同一个态射 \(x\rightarrow y\) 决定，因为前面的边全是恒等；故它是集合 \(\operatorname{Hom}_{\mathcal A}(x,y)\) 对应的离散单纯集。

\ProseParagraphBreak
把映射空间只取连通分支，就得到同伦范畴 \(h\mathcal C\)：对象仍为顶点，态射集合为
\[
 \operatorname{Hom}_{h\mathcal C}(x,y)
 =\pi_0\operatorname{Map}_{\mathcal C}(x,y).
\]
这里 \(\pi_0\) 是由一维单纯形的两个端点生成的等价关系的商集，恒等态射为 \([s_0x]\)。对边 \(f:x\rightarrow y\)、\(g:y\rightarrow z\)，复合 \([g][f]\) 由填充 \(\Lambda_1^2\) 后的第三条边定义；填充与代表的选择在 \(\pi_0\) 上给出同一结果，结合律由三维填充保证。左右同伦的比较、复合的良定义及范畴公理见\cite[Remark 1.2.3.6; Propositions 1.2.3.7--1.2.3.9]{Lurie2009HigherToposTheory}。在普通范畴神经的情形，这恢复 \(hN(\mathcal A)\cong\mathcal A\)。

\ProseParagraphBreak
复形的例子则有非离散的映射空间。对域 \(k\) 上的复形，\subsecref{b4:subsec:B04:02}给出的稳定增强满足
\[
 \pi_i\operatorname{Map}(X,Y)
 \cong H^{-i}\operatorname{Hom}^{\bullet}_k(X,Y)
 \qquad(i\ge0),
\]
正次同伦群取零映射为基点。例如 \(X=k\)、\(Y=k[1]\) 时，映射空间只有一个连通分支，却有 \(\pi_1\cong k\)。本节的角模型说明映射空间应具备什么相容性；从 DG 结构得到这一模型还需要单纯化与 DG 神经的构造及比较，它们的范围和次数约定见该小节。
